\chapter{总结与展望}
本文围绕三维光场人脸光照编辑这一核心课题，针对静态与动态场景下的技术痛点，构建了从基础方法到优化框架的完整技术体系，通过理论创新与实验验证，有效解决了现有方法中存在的采集设备复杂、多视角一致性不足、动态时序闪烁等关键问题，主要研究成果与核心内容如下。
\section{研究内容与创新}
提出了一种关键信息驱动的三维光场静态人脸光照编辑方法。针对静态三维光场人像光照编辑中 “多视图独立编辑导致一致性差” 与 “密集采集设备成本高昂” 的双重痛点，提出了关键信息驱动的光照编辑方法，核心创新与成果如下：构建了 “高光去除 - 光照编辑 - 密集视点生成” 三级框架：通过双重检测标准（亮度阈值 + 漫反射分量校验）精准定位高光区域，结合动态窗口自适应替换算法消除高光干扰，为后续光照编辑提供干净的输入基础；采用双路径光照编辑策略，既支持从 HDR 图像提取 SH 系数实现环境光重光照，又能通过 BERT 模型解析自然语言指令，实现光照方向、强度的语义化调控；基于 NeRF 体渲染框架，生成水平 ±30° 范围内 90 个视角的密集视差图，满足光场显示需求。

实验验证成效显著：客观实验中，与 NVPR、LRPI 等主流方法相比，本文方法在 FID（33.55 vs 36.54/36.15）、PSNR（32.06 vs 28.66/28.44）、SSIM（0.8954 vs 0.6288/0.8145）及 LPIPS（0.19 vs 0.30/0.23）等指标上均实现最优，光照误差（LE=0.0607）低于对比方法；单帧推理效率大幅提升，在 RTX 3090 GPU 上处理 1K 分辨率图像仅需 25.92s，远快于 NVPR（230.45s）与 LRPI（180.21s）；主观实验中，在三维光场显示器上的立体感知得票率达 89.19\%，多视角光照一致性得票率 92.29\%，显著优于单视图逐帧编辑方法，证明该方法在保证光照一致性的前提下，实现了简单高效的静态光场人像光照编辑。


提出了一种空间双维联合优化的动态光照编辑框架针。对静态方法扩展至视频序列时出现的时序闪烁与多视角光照失配问题，提出动态光照一致性优化框架，核心创新与成果如下：
时序一致性优化模块：采用共享编码器提取相邻帧稳定语义特征，通过光流估计网络实现跨帧空间对齐，将时间一致性问题转化为空间对齐问题；设计光照分量提取网络解耦几何与光照特征，仅对光照响应分量施加 “特征 L1 + 像素残差 L2” 双重损失约束，在抑制帧间闪烁的同时保留表情与姿态的自然动态。实验表明，该模块使时序稳定性指标 WE 降至 0.0004，帧间 LPIPS 低至 0.1912，显著优于 SMFR（WE=0.0022）与 E4E（帧间 LPIPS=0.2344）。
多视角一致性优化模块：基于相机内外参构建投影矩阵，结合 NeRF 获得统一三维几何模型，通过三维表面采样将多视角光照约束转化为三维点光照观测问题；引入多因子加权聚合机制（可见性 + 法线夹角 + 光照置信度），实现可靠的光照统一估计，再通过一致性回投影优化还原至二维视角。该模块使多视角光照误差 LE 降至 0.7922，光照不一致度 LI 达 0.2644，均优于对比方法。

综合性能验证：客观实验中，本文方法在 FID（33.5501）、PSNR（32.0226）等生成质量指标上保持领先，身份一致性 ID 达 0.5122，优于对比方法；主观实验中，动态光照连续性、多视角一致性及整体三维沉浸感综合评分达 4.73，较 PTI 方法（3.08）提升 53.57；消融实验验证了时序一致性模块儿的重要性。

\section{未来与展望}
本文虽在三维光场人脸光照编辑领域取得了阶段性成果，但仍存在可拓展与优化的方向，结合技术发展趋势与实际应用需求，未来研究可聚焦以下方面：

场景通用性拓展：当前方法主要针对人像场景设计，后续可研究非人像场景（如工业零件、自然景物）的几何与光照解耦方法，通过优化三维重建与光照估计策略，提升算法对复杂轮廓、多样材质的适应性，拓宽技术的应用场景。

实时性优化与硬件适配：现有方法在处理高分辨率视频时，实时推理性能仍有提升空间。未来可通过模型轻量化（如剪枝、量化、蒸馏）、硬件加速（如 GPU/TPU 专用算子优化）及算法简化（如高效光流估计、稀疏采样策略），降低计算复杂度，实现高清三维光场视频的实时光照编辑，满足直播、AR/VR 等实时交互场景的需求。

生成式 AI 融合与精细编辑：结合生成式对抗网络（GAN）、扩散模型等先进生成式 AI 技术，实现更精细的光照细节编辑（如局部光斑调控、纹理光影交互）；探索多模态输入（文本、语音、草图）的光照编辑方式，进一步降低用户操作门槛，提升交互灵活性与编辑自由度。

多场景光照交互与动态适配：当前方法主要针对固定光照编辑需求，未来可研究动态场景下的光照自适应调整策略，实现光照与场景运动（如人脸大幅度动作、视角剧烈切换）的实时适配；探索多物体光场场景中的光照交互机制，支持不同物体间的光照反射、阴影投射等物理效果模拟，提升光场显示的整体真实感与沉浸感。

